%-*-latex-*-
%-*-latex-*-
\newcommand\COURSE{ciss240}
\newcommand\ASSESSMENT{q1202}
\newcommand\ASSESSMENTTYPE{Quiz}
\newcommand\POINTS{\textwhite{xxx/xxx}}

\input{myquizpreamble}
\input{yliow}
\input{\COURSE}
%-*-latex-*-

%-*-latex-*-

%-*-latex-*-
\renewcommand\TITLE{\ASSESSMENTTYPE \ \ASSESSMENT}

\renewcommand\EMAIL{}
\renewcommand\AUTHOR{eagraham1@cougars.ccis.edu}

\textwidth=6in
\begin{document}
\topmatterthree

%------------------------------------------------------------------------------
\nextq
Write a program that continually get an integer from the
user prints \verb!zero! if the input is \verb!0!,
prints \verb!unit! is the input is \verb!1!,
otherwise it prints \verb!prime! if the input is a prime number
or it prints \verb!composite! if the input is not prime.
If the user enters a negative number the program prints \verb!error!.
To stop the loop and ends the program, the user enters \verb!-1!.
Here's a sample test run:
\begin{console}[commandchars=\\\{\},fontsize=\footnotesize]
\userinput{0}
zero
\userinput{1}
unit
\userinput{2}
prime
\userinput{3}
prime
\userinput{4}
composite
\userinput{-5}
error
\userinput{-1}
\end{console}
\ANSWER
\begin{answercode}
#include <iostream>
#include <cmath>

bool is_prime(int num)
{
    if (num < 2) return false;          
    for (int i = 2; i <= std::sqrt(num); ++i)
    {
        if (num % i == 0) return false;
    }
    return true; 
}

int main()
{
    int input;

    while (true)
    {
        std::cin >> input; 

        if (input == -1)
        {              
            break;
        }
        else if (input == 0)
        {
            std::cout << "zero\n";
        }
        else if (input == 1)
        {
            std::cout << "unit\n";
        }
        else if (input < 0)
        {
            std::cout << "error\n";
        }
        else if (is_prime(input))
        {
            std::cout << "prime\n";
        }
        else
        {
            std::cout << "composite\n";
        }
    }

    return 0;
}
\end{answercode}

%------------------------------------------------------------------------------
%-*-latex-*-
\newpage
\input{instructions}
\end{document}

